% -------------------------- 矢量分析（教材 附录一 / 课件 第0章） --------------------------
% 说明：本文件为子文件，由主文档 \input 引入，不含 document 环境与重复宏包
% 依赖宏包（主文档 preamble 已装）：tcolorbox、adjustbox、booktabs、amsmath
%
% 彩色公式框：本文件定义一次，供 4.静电场.tex 等同目录子文件复用（勿重复定义）
\newenvironment{formulaBox}{%
    \par\begin{adjustbox}{max width=\linewidth}%
    \begin{tcolorbox}[colback=blue!8!white,colframe=transparent,boxrule=0pt,halign=center,boxsep=2pt,left=4pt,right=4pt,top=3pt,bottom=3pt]%
}{%
    \end{tcolorbox}\end{adjustbox}\par
}

% -------------------------- 正文 --------------------------
\chapter{矢量分析}

矢量分析是研究电磁场理论的重要数学工具。它对简化公式、明确概念、掌握观察规律的实质有极大裨益，
也是学习电磁场入门的数学基础。本章（教材列为附录一，课件称“第 0 章”）只讨论梯度、散度、旋度的概念、物理意义
与直角坐标系下的计算，并给出三种正交坐标系中的常用公式；不追求纯数学论证，侧重阐明物理概念与
应用手段。

\section{正交坐标系}

\subsection{场的描述与坐标系的选择}

电磁场的场量是空间点与时间的函数。为了同时刻画“大小”“方向”和“变化率”，必须选一套坐标基。
当坐标系的三个坐标基矢量互相垂直、且都取单位矢量时，称为\textbf{正交坐标系}，它同时满足
正交性与归一性：
\begin{formulaBox}
\[
\vec{e}_i\cdot\vec{e}_j=\delta_{ij}=
\begin{cases}1, & i=j\\ 0, & i\neq j\end{cases}
\qquad |\vec{e}_i|=1
\]
\end{formulaBox}

题目的边界形状决定坐标系的选择——边界若能与某坐标系的坐标面相重合，问题往往能大幅简化。
这正是后面“分离变量法”能否奏效的前提。

\subsection{三种常用正交坐标系}

\begin{table}[htbp]
\centering
{\small\textbf{表 1-1\quad 三种正交坐标系的坐标、基矢量与适用场域}}\\[2pt]
\small
\begin{tabular}{@{}llll@{}}
\toprule
坐标系 & 坐标 & 单位矢量 & 基矢量是否随位置变化 \\
\midrule
直角（笛卡儿） & $(x,y,z)$ & $\vec{e}_x,\vec{e}_y,\vec{e}_z$ & 常矢量，不随位置变化 \\
圆柱 & $(\rho,\varphi,z)$ & $\vec{e}_\rho,\vec{e}_\varphi,\vec{e}_z$ & $\vec{e}_\rho,\vec{e}_\varphi$ 随 $\varphi$ 变 \\
球 & $(r,\theta,\varphi)$ & $\vec{e}_r,\vec{e}_\theta,\vec{e}_\varphi$ & 三者均随位置变化 \\
\bottomrule
\end{tabular}
\end{table}

名称约定：$\rho$ 为圆柱坐标的径向距离，$\varphi$ 为方位角；$r$ 为球坐标的矢径长度，
$\theta$ 为极角（与 $z$ 轴夹角），$\varphi$ 为方位角。三者取值范围
$\rho\ge 0$、$r\ge 0$、$0\le\theta\le\pi$、$0\le\varphi<2\pi$。

基矢量是否随位置变化，决定了求导与 $\nabla$ 表达式的繁简：直角坐标的三个基矢量方向固定，
故可直接对分量求导，运算最简；圆柱坐标的 $\vec{e}_\rho$、$\vec{e}_\varphi$ 随 $\varphi$ 改变，
是“变矢量”，对它们求导不为零，这正是圆柱坐标下 $\nabla$ 表达式出现 $1/\rho$ 因子的根源；
球坐标的三个基矢量都随位置改变。相应地，直角坐标适用于无限大平面、矩形边界、平板电极；
圆柱坐标适用于同轴电缆、长直导线；球坐标适用于点电荷、同心球壳。

\subsection{坐标变换与单位矢量的变换}

同一场点在不同坐标系中的坐标之间可以互换。设 $x,y,z$ 为直角坐标，则
\begin{formulaBox}
\[
\begin{aligned}
&\text{直角}\to\text{圆柱}:\quad \rho=\sqrt{x^2+y^2},\quad \varphi=\arctan\frac{y}{x},\quad z=z\\
&\text{圆柱}\to\text{直角}:\quad x=\rho\cos\varphi,\quad y=\rho\sin\varphi,\quad z=z\\
&\text{直角}\to\text{球}:\quad r=\sqrt{x^2+y^2+z^2},\quad \theta=\arccos\frac{z}{r},\quad \varphi=\arctan\frac{y}{x}\\
&\text{球}\to\text{直角}:\quad x=r\sin\theta\cos\varphi,\quad y=r\sin\theta\sin\varphi,\quad z=r\cos\theta
\end{aligned}
\]
\end{formulaBox}

反三角函数的象限由 $x$、$y$ 的符号共同确定（只用 $\arctan(y/x)$ 会丢失象限信息，这是坐标变换
最常见的错误）。单位矢量之间的变换为
\begin{formulaBox}
\[
\begin{aligned}
\vec{e}_\rho&=\cos\varphi\,\vec{e}_x+\sin\varphi\,\vec{e}_y, &
\vec{e}_\varphi&=-\sin\varphi\,\vec{e}_x+\cos\varphi\,\vec{e}_y\\[2pt]
\vec{e}_r&=\sin\theta\cos\varphi\,\vec{e}_x+\sin\theta\sin\varphi\,\vec{e}_y+\cos\theta\,\vec{e}_z, &
\vec{e}_\theta&=\cos\theta\cos\varphi\,\vec{e}_x+\cos\theta\sin\varphi\,\vec{e}_y-\sin\theta\,\vec{e}_z
\end{aligned}
\]
\end{formulaBox}

线元、面元、体元是推导管积分的基础，必须记准：
\begin{formulaBox}
\[
\begin{aligned}
\text{直角}:\quad &\mathrm{d}l=\mathrm{d}x\,\vec{e}_x+\mathrm{d}y\,\vec{e}_y+\mathrm{d}z\,\vec{e}_z,\quad \mathrm{d}V=\mathrm{d}x\,\mathrm{d}y\,\mathrm{d}z\\
\text{圆柱}:\quad &\mathrm{d}l=\mathrm{d}\rho\,\vec{e}_\rho+\rho\,\mathrm{d}\varphi\,\vec{e}_\varphi+\mathrm{d}z\,\vec{e}_z,\quad \mathrm{d}V=\rho\,\mathrm{d}\rho\,\mathrm{d}\varphi\,\mathrm{d}z\\
\text{球}:\quad &\mathrm{d}l=\mathrm{d}r\,\vec{e}_r+r\,\mathrm{d}\theta\,\vec{e}_\theta+r\sin\theta\,\mathrm{d}\varphi\,\vec{e}_\varphi,\quad \mathrm{d}V=r^2\sin\theta\,\mathrm{d}r\,\mathrm{d}\theta\,\mathrm{d}\varphi
\end{aligned}
\]
\end{formulaBox}

\subsection{源点、场点与距离矢量}

计算场量时要区分两个点：产生场的电荷所在处称\textbf{源点}，坐标带撇
（$x',y',z'$ 或 $\vec{r}'$）；待求场量的点称\textbf{场点}，坐标不带撇（$x,y,z$ 或 $\vec{r}$）。
两者之间的\textbf{距离矢量}规定由源点指向场点：
\begin{formulaBox}
\[
\vec{R}=\vec{r}-\vec{r}'=R\,\vec{e}_R,\qquad R=|\vec{r}-\vec{r}'|=\sqrt{(x-x')^2+(y-y')^2+(z-z')^2}
\]
\end{formulaBox}

由此导出电磁场中使用频率最高的两个微分关系（$\nabla$ 对场点作用，$\nabla'$ 对源点作用）：
\begin{formulaBox}
\[
\nabla R=\frac{\vec{R}}{R}=\vec{e}_R,\qquad
\nabla'\left(\frac{1}{R}\right)=-\nabla\left(\frac{1}{R}\right)=\frac{\vec{R}}{R^3}=\frac{\vec{e}_R}{R^2}
\]
\end{formulaBox}

\corepoint{注意 $\nabla$ 与 $\nabla'$ 只差一个负号，但作用对象完全不同：前者对场点求导，后者对源点求导。
写积分式时若把 $\nabla$ 提到积分号外，必须确认被积函数中只有源点坐标与积分变量有关。}

\section{矢量代数}

\subsection{标量与矢量}

\textbf{标量}：坐标轴转动时数值不变的物理量，如质量、温度、电位、能量、电阻、电荷量。
\textbf{矢量}：坐标轴转动时按转动矩阵变换的物理量，如位移、速度、力、电场强度、电流密度。
矢量在直角坐标系中写成三个分量与基矢量之积的和，其模为
\begin{formulaBox}
\[
\vec{A}=A_x\vec{e}_x+A_y\vec{e}_y+A_z\vec{e}_z,\qquad
|\vec{A}|=\sqrt{A_x^2+A_y^2+A_z^2},\qquad
\vec{e}_A=\frac{\vec{A}}{|\vec{A}|}
\]
\end{formulaBox}

\subsection{矢量点乘（标量积）}

定义为两矢量模与夹角余弦之积，结果是一个标量：
\begin{formulaBox}
\[
\vec{A}\cdot\vec{B}=|\vec{A}||\vec{B}|\cos(\vec{A},\vec{B})
=A_xB_x+A_yB_y+A_zB_z
\]
\end{formulaBox}

\textbf{理解要点}：点乘等于一个矢量在另一矢量方向上的投影与另一矢量模的乘积——即“$\vec{A}$ 在
$\vec{B}$ 上的投影乘以 $|\vec{B}|$”。因此

\begin{itemize}
    \item $\vec{A}\cdot\vec{B}=0\iff \vec{A}\perp\vec{B}$（正交判据）；
    \item 满足交换律 $\vec{A}\cdot\vec{B}=\vec{B}\cdot\vec{A}$ 与分配律
    $\vec{A}\cdot(\vec{B}+\vec{C})=\vec{A}\cdot\vec{B}+\vec{A}\cdot\vec{C}$；
    \item 物理原型：力做功 $W=\vec{F}\cdot\vec{l}$，电通量 $\Phi=\vec{D}\cdot\vec{S}$。
\end{itemize}

\subsection{矢量叉乘（矢量积）}

结果仍为矢量，大小等于两矢量张成的平行四边形面积，方向由右手螺旋法则确定：
\begin{formulaBox}
\[
\vec{A}\times\vec{B}=
\begin{vmatrix}
\vec{e}_x & \vec{e}_y & \vec{e}_z\\
A_x & A_y & A_z\\
B_x & B_y & B_z
\end{vmatrix}
=|\vec{A}||\vec{B}|\sin(\vec{A},\vec{B})\,\vec{e}_n
\]
\end{formulaBox}

\begin{itemize}
    \item $\vec{A}\times\vec{B}=0\iff \vec{A}\parallel\vec{B}$（平行判据）；
    \item 反交换律 $\vec{A}\times\vec{B}=-\vec{B}\times\vec{A}$，且不满足结合律；
    \item 物理原型：力矩 $\vec{M}=\vec{r}\times\vec{F}$，洛伦兹力 $\vec{F}=q\vec{v}\times\vec{B}$，
    以及“绕行方向与曲面法向成右手螺旋”的有向面积规定。
\end{itemize}

展开行列式时\textbf{$y$ 分量带负号}，这是叉乘最容易出错的一处：
\begin{formulaBox}
\[
\vec{A}\times\vec{B}=(A_yB_z-A_zB_y)\vec{e}_x+(A_zB_x-A_xB_z)\vec{e}_y+(A_xB_y-A_yB_x)\vec{e}_z
\]
\end{formulaBox}

\subsection{混合积与矢量的分解}

三个矢量的混合积 $\vec{A}\cdot(\vec{B}\times\vec{C})$ 是一个标量，其绝对值等于三矢量张成的
平行六面体体积，且满足轮换对称性；二重叉乘则服从“展开法则”：
\begin{formulaBox}
\[
\vec{A}\cdot(\vec{B}\times\vec{C})=\vec{B}\cdot(\vec{C}\times\vec{A})=\vec{C}\cdot(\vec{A}\times\vec{B}),\qquad
\vec{A}\times(\vec{B}\times\vec{C})=(\vec{A}\cdot\vec{C})\vec{B}-(\vec{A}\cdot\vec{B})\vec{C}
\]
\end{formulaBox}

利用正交基，任一矢量都可沿三个互相垂直的方向唯一分解，这正是把矢量方程化为三个标量方程的依据。

\section{标量场与矢量场}

\subsection{场的概念}

分布着某种物理量的空间区域称为该物理量的\textbf{场}。若物理量是标量，称为\textbf{标量场}，
记作 $f(x,y,z;t)$，如温度场、电位场；若是矢量，称为\textbf{矢量场}，记作
$\vec{A}(x,y,z;t)$，如速度场、电场、磁场。不随时间变化的场称\textbf{静态场}，
否则称\textbf{时变场}（动态场）。

\subsection{标量场的等值面与等值线}

等值面是函数值相等的空间曲面，其方程为
\begin{formulaBox}
\[
f(x,y,z)=C
\]
\end{formulaBox}

取不同的 $C$ 得到一族等值面，它们充满整个场域。由于 $f$ 一般是单值函数，\textbf{等值面互不相交}。
电位场中的等位面、温度场中的等温面都是等值面；平行平面场（二维场）中等值面退化为\textbf{等值线}。

等值面的疏密可直接读出场的强弱：相邻等值面间差值相同时，面愈密处场量变化愈剧烈，梯度愈大。

\subsection{矢量场的矢量线}

矢量线是这样的曲线：线上每一点的切线方向与该点矢量场的方向相同，如静电场中的电力线、
流速场中的流线。矢量线充满场域但互不相交（否则该点场方向不唯一）。

因为线元 $\mathrm{d}\vec{l}$ 与 $\vec{A}$ 同向，故矢量线的微分方程（亦称矢量线方程）为
\begin{formulaBox}
\[
\vec{A}\times\mathrm{d}\vec{l}=0
\quad\Longleftrightarrow\quad
\frac{\mathrm{d}x}{A_x}=\frac{\mathrm{d}y}{A_y}=\frac{\mathrm{d}z}{A_z}
\]
\end{formulaBox}

\corepoint{矢量线方程是“成比例”的微分方程组，求解时任取其中一对作常微分方程积分，得到两个积分曲面，
其交线即为矢量线。若给定过某点，用该点坐标定出积分常数。}

作过某曲面 $S$ 上全部点的矢量线，这些线围成的管状区域称\textbf{矢量管}。矢量线只表征方向；
若要同时表征大小，则规定\textbf{穿过单位垂直面积的矢量线条数正比于场的大小}——线愈密处场愈强。

\section{标量场的梯度}

\subsection{方向导数}

标量场 $f(P)$ 从给定点沿某方向 $l$ 的变化率称为\textbf{方向导数}，记作 $\partial f/\partial l$。
设 $l$ 的方向余弦为 $\cos\alpha,\cos\beta,\cos\gamma$，$\vec{e}_l$ 为 $l$ 方向的单位矢量，则
\begin{formulaBox}
\[
\frac{\partial f}{\partial l}=\frac{\partial f}{\partial x}\cos\alpha
+\frac{\partial f}{\partial y}\cos\beta
+\frac{\partial f}{\partial z}\cos\gamma,
\qquad
\vec{e}_l=\vec{e}_x\cos\alpha+\vec{e}_y\cos\beta+\vec{e}_z\cos\gamma
\]
\end{formulaBox}

方向导数大于零表示沿该方向 $f$ 增大，小于零表示减小，等于零表示沿该方向 $f$ 不变
（即沿等值面的切向）。它是一个\textbf{与方向有关}的量：同一点沿不同方向，变化率不同。

\subsection{梯度的定义、大小与方向}

既然从一点出发有无穷多个方向，自然要问：哪个方向变化率最大？最大值是多少？把上式的三个偏导数
记成一个矢量 $\vec{g}$，则方向导数就是 $\vec{g}$ 在 $\vec{e}_l$ 上的投影：
\begin{formulaBox}
\[
\frac{\partial f}{\partial l}=\vec{g}\cdot\vec{e}_l=|\vec{g}|\cos(\vec{g},\vec{e}_l)
\]
\end{formulaBox}

当 $\vec{e}_l$ 与 $\vec{g}$ 同向时方向导数取最大值 $|\vec{g}|$，反向时取最小值 $-|\vec{g}|$，
垂直时为零。于是定义该矢量为标量场的\textbf{梯度}：
\begin{formulaBox}
\[
\operatorname{grad} f=\nabla f
=\frac{\partial f}{\partial x}\vec{e}_x+\frac{\partial f}{\partial y}\vec{e}_y+\frac{\partial f}{\partial z}\vec{e}_z
\]
\end{formulaBox}

\begin{itemize}
    \item 梯度是\textbf{矢量}，且是空间坐标点的函数；其\textbf{方向}为该点最大变化率的方向，
    其\textbf{模}等于该点的最大方向导数。
    \item 已知梯度即可求任一方向的方向导数：$\partial f/\partial l=\nabla f\cdot\vec{e}_l$。
    \item 梯度\textbf{垂直于等值面}，且指向函数值增大的方向；等值面法向单位矢量
    $\vec{e}_n=\pm\nabla f/|\nabla f|$。
    \item $\nabla f=0$ 的点是 $f$ 的驻点；若在场域内处处 $\nabla f=0$，则 $f$ 为常数。
\end{itemize}

在圆柱、球坐标系中，梯度的表达式多出 $1/\rho$、$1/r$ 因子：
\begin{formulaBox}
\[
\nabla f=\frac{\partial f}{\partial\rho}\vec{e}_\rho+\frac{1}{\rho}\frac{\partial f}{\partial\varphi}\vec{e}_\varphi+\frac{\partial f}{\partial z}\vec{e}_z,
\qquad
\nabla f=\frac{\partial f}{\partial r}\vec{e}_r+\frac{1}{r}\frac{\partial f}{\partial\theta}\vec{e}_\theta+\frac{1}{r\sin\theta}\frac{\partial f}{\partial\varphi}\vec{e}_\varphi
\]
\end{formulaBox}

\subsection{哈密顿算子}

为书写与运算方便，引入\textbf{矢量性微分算子}（哈密顿算子，读作 nabla）：
\begin{formulaBox}
\[
\nabla=\vec{e}_x\frac{\partial}{\partial x}+\vec{e}_y\frac{\partial}{\partial y}+\vec{e}_z\frac{\partial}{\partial z}
\]
\end{formulaBox}

它同时具备两种性质：既有\textbf{微分性}（对跟随其后的函数求导），又有\textbf{矢量性}
（可像矢量一样作点乘、叉乘，但不能与之交换顺序）。因此同一符号可以表示三种运算：
\begin{formulaBox}
\[
\begin{aligned}
\nabla f&:\quad \text{梯度（标量场}\to\text{矢量场）}\\
\nabla\cdot\vec{A}&:\quad \text{散度（矢量场}\to\text{标量场）}\\
\nabla\times\vec{A}&:\quad \text{旋度（矢量场}\to\text{矢量场）}
\end{aligned}
\]
\end{formulaBox}

算子作用于不同对象时遵守普通求导法则，例如对两个标量场之积：
\begin{formulaBox}
\[
\nabla(au+bv)=a\nabla u+b\nabla v,\qquad \nabla(uv)=u\nabla v+v\nabla u
\]
\end{formulaBox}

（$a,b$ 为常数。）注意 $\nabla$ 不能写成 $\vec{A}\cdot\nabla$ 与 $\nabla\cdot\vec{A}$ 混用：
前者是标量算子，后者是散度，两者含义不同。

\subsection{拉普拉斯算子}

两个 $\nabla$ 的标积是一个标量算子，称\textbf{拉普拉斯算子}：
\begin{formulaBox}
\[
\nabla^2=\nabla\cdot\nabla=\frac{\partial^2}{\partial x^2}+\frac{\partial^2}{\partial y^2}+\frac{\partial^2}{\partial z^2}
\]
\end{formulaBox}

它作用在标量函数上得标量（拉普拉斯运算），作用在矢量函数上则对三个分量分别运算：
\begin{formulaBox}
\[
\nabla^2 f=\frac{\partial^2 f}{\partial x^2}+\frac{\partial^2 f}{\partial y^2}+\frac{\partial^2 f}{\partial z^2},
\qquad
\nabla^2\vec{A}=(\nabla^2A_x)\vec{e}_x+(\nabla^2A_y)\vec{e}_y+(\nabla^2A_z)\vec{e}_z
\]
\end{formulaBox}

圆柱、球坐标系中的标量拉普拉斯算子为
\begin{formulaBox}
\[
\nabla^2 f=\frac{1}{\rho}\frac{\partial}{\partial\rho}\left(\rho\frac{\partial f}{\partial\rho}\right)
+\frac{1}{\rho^2}\frac{\partial^2 f}{\partial\varphi^2}+\frac{\partial^2 f}{\partial z^2}
\]
\end{formulaBox}
\begin{formulaBox}
\[
\nabla^2 f=\frac{1}{r^2}\frac{\partial}{\partial r}\left(r^2\frac{\partial f}{\partial r}\right)
+\frac{1}{r^2\sin\theta}\frac{\partial}{\partial\theta}\left(\sin\theta\frac{\partial f}{\partial\theta}\right)
+\frac{1}{r^2\sin^2\theta}\frac{\partial^2 f}{\partial\varphi^2}
\]
\end{formulaBox}

\corepoint{圆柱坐标下 $\nabla^2 f$ 的第一项必须写成 $\frac{1}{\rho}\frac{\partial}{\partial\rho}\left(\rho\frac{\partial f}{\partial\rho}\right)$，
不能写成 $\frac{\partial^2 f}{\partial\rho^2}$；球坐标下也不等于 $\frac{\partial^2 f}{\partial r^2}$。
少了 $\rho$、$r^2$ 因子是解泊松/拉普拉斯方程时最常见的错误。}

\section{矢量场的通量与散度}

\subsection{通量}

在矢量场中取一个有向曲面 $S$（$\vec{e}_n$ 为其单位法向），矢量 $\vec{A}$ 在 $S$ 上的面积分
称为 $\vec{A}$ 穿过 $S$ 的\textbf{通量}：
\begin{formulaBox}
\[
\Phi=\iint_S\vec{A}\cdot\mathrm{d}\vec{S}=\iint_S\vec{A}\cdot\vec{e}_n\,\mathrm{d}S
=\iint_S\left(A_x\,\mathrm{d}y\mathrm{d}z+A_y\,\mathrm{d}z\mathrm{d}x+A_z\,\mathrm{d}x\mathrm{d}y\right)
\]
\end{formulaBox}

通量的数学意义是\textbf{穿过有向曲面的矢量线条数}；物理意义随场而定：流速场的通量是单位时间
穿过 $S$ 的体积流量，电流密度 $\vec{J}$ 的通量是电流，电位移 $\vec{D}$ 的通量是自由电荷量。

对\textbf{闭合曲面}，习惯取外法向为正，此时通量表示穿出的净通量，可据其正负判断面内源的性质：
\begin{formulaBox}
\[
\Phi=\oiint_S\vec{A}\cdot\mathrm{d}\vec{S}
=\begin{cases}
>0, & \text{面内有正源（发出矢量线）}\\
<0, & \text{面内有负源（汇集矢量线）}\\
=0, & \text{无源，或正源与负源相抵}
\end{cases}
\]
\end{formulaBox}

通量是积分量，只反映曲面所围区域内的\textbf{整体}情况，说不清面内某一点的性质——因此引出散度。

\subsection{散度的定义与计算}

散度是“通量的体密度”，即单位体积发出的通量：
\begin{formulaBox}
\[
\operatorname{div}\vec{A}=\nabla\cdot\vec{A}
=\lim_{\Delta V\to 0}\frac{1}{\Delta V}\oiint_S\vec{A}\cdot\mathrm{d}\vec{S}
\]
\end{formulaBox}

在直角坐标系中可展开为三个偏导数之和：
\begin{formulaBox}
\[
\nabla\cdot\vec{A}=\frac{\partial A_x}{\partial x}+\frac{\partial A_y}{\partial y}+\frac{\partial A_z}{\partial z}
\]
\end{formulaBox}

柱坐标与球坐标下的散度为
\begin{formulaBox}
\[
\nabla\cdot\vec{A}=\frac{1}{\rho}\frac{\partial(\rho A_\rho)}{\partial\rho}
+\frac{1}{\rho}\frac{\partial A_\varphi}{\partial\varphi}+\frac{\partial A_z}{\partial z}
\]
\end{formulaBox}
\begin{formulaBox}
\[
\nabla\cdot\vec{A}=\frac{1}{r^2}\frac{\partial(r^2A_r)}{\partial r}
+\frac{1}{r\sin\theta}\frac{\partial(\sin\theta A_\theta)}{\partial\theta}
+\frac{1}{r\sin\theta}\frac{\partial A_\varphi}{\partial\varphi}
\]
\end{formulaBox}

\begin{itemize}
    \item 散度是\textbf{标量}，是空间坐标点的函数，描述该点通量源的强度与正负。
    \item $\nabla\cdot\vec{A}>0$：该点有发出通量线的正源；$<0$：有汇集通量线的负源；
    $=0$：该点无源。
    \item 若场中处处 $\nabla\cdot\vec{A}=0$，称\textbf{无源场}（如恒定磁场）。
    \item 静电场中 $\nabla\cdot\vec{D}=\rho_f$：静止电荷就是电场的散度源。
\end{itemize}

\subsection{高斯散度定理}

把散度的定义对体积积分，就得到电磁场理论中使用最广的积分变换定理：
\begin{formulaBox}
\[
\oiint_S\vec{A}\cdot\mathrm{d}\vec{S}=\iiint_V\nabla\cdot\vec{A}\,\mathrm{d}V
\qquad (S \text{ 为包围 } V \text{ 的闭合曲面})
\]
\end{formulaBox}

它把\textbf{闭合曲面上的面积分}与\textbf{体积分}互相转化，是“微分形式 $\Leftrightarrow$ 积分形式”场方程的桥梁：
静电场的高斯定理、恒定磁场的磁通连续性方程、时变场的电荷守恒方程都由它沟通。

\corepoint{应用散度定理前须检查两件事：（1） $S$ 必须是闭合面且取外法向；（2） $\vec{A}$ 在 $V$ 内须有连续的一阶偏导数。
若场域内有奇异点（如点电荷所在处），必须把这些点挖去，用“挖补法”在多连通域上分区应用。}

\section{矢量场的环量与旋度}

\subsection{环量}

矢量场沿有向闭合曲线的线积分称为\textbf{环量}：
\begin{formulaBox}
\[
\Gamma=\oint_L\vec{A}\cdot\mathrm{d}\vec{l}
\]
\end{formulaBox}

环量的物理意义随场而定：力场中的环量是力沿闭合路径做的功；流速场中的环量反映流体是否作涡旋运动。
$\Gamma=0$ 表示沿该回路无涡旋源；$\Gamma\neq 0$ 表示回路包围的区域内必有产生涡旋的源。
但环量和通量一样是积分量，说不清回路内每一点的涡旋强度——因此引出旋度。

\subsection{环量面密度与旋度}

围绕一点可取无穷多个方向的面元 $S$，每个面元对应一个环量面密度 $\Gamma/\Delta S$。其中最大者
的方向与大小就定义了旋度。在直角坐标系中：
\begin{formulaBox}
\[
\operatorname{rot}\vec{A}=\nabla\times\vec{A}=
\begin{vmatrix}
\vec{e}_x & \vec{e}_y & \vec{e}_z\\[2pt]
\dfrac{\partial}{\partial x} & \dfrac{\partial}{\partial y} & \dfrac{\partial}{\partial z}\\[6pt]
A_x & A_y & A_z
\end{vmatrix}
\]
\end{formulaBox}

\begin{itemize}
    \item 旋度是\textbf{矢量}：其\textbf{方向}为使环量面密度取最大值时面元法线的方向，
    其\textbf{模}等于该点环量面密度的最大值。
    \item $\nabla\times\vec{A}\neq 0$：该点矢量线有旋（存在涡旋源）；$=0$：该点无旋。
    若场中处处 $\nabla\times\vec{A}=0$，称\textbf{无旋场}（如静电场）。
    \item 若 $\nabla\times\vec{A}=\vec{J}\neq 0$，称 $\vec{J}$ 为\textbf{旋度源}（涡旋源）。
\end{itemize}

柱坐标与球坐标下的旋度为（注意行列式中的 $\rho$、$r$、$\sin\theta$ 因子）
\begin{formulaBox}
\[
\nabla\times\vec{A}=\frac{1}{\rho}
\begin{vmatrix}
\vec{e}_\rho & \rho\vec{e}_\varphi & \vec{e}_z\\[2pt]
\dfrac{\partial}{\partial\rho} & \dfrac{\partial}{\partial\varphi} & \dfrac{\partial}{\partial z}\\[6pt]
A_\rho & \rho A_\varphi & A_z
\end{vmatrix}
\]
\end{formulaBox}
\begin{formulaBox}
\[
\nabla\times\vec{A}=\frac{1}{r^2\sin\theta}
\begin{vmatrix}
\vec{e}_r & r\vec{e}_\theta & r\sin\theta\,\vec{e}_\varphi\\[2pt]
\dfrac{\partial}{\partial r} & \dfrac{\partial}{\partial\theta} & \dfrac{\partial}{\partial\varphi}\\[6pt]
A_r & rA_\theta & r\sin\theta\,A_\varphi
\end{vmatrix}
\]
\end{formulaBox}

\subsection{斯托克斯定理}

把旋度的定义沿曲面叠加，得到另一个基本积分变换定理：
\begin{formulaBox}
\[
\oint_L\vec{A}\cdot\mathrm{d}\vec{l}=\iint_S(\nabla\times\vec{A})\cdot\mathrm{d}\vec{S}
\]
\end{formulaBox}

其中 $S$ 是以闭合曲线 $L$ 为周界的曲面，$L$ 的绕行方向与 $S$ 的法向成\textbf{右手螺旋}关系。
它把\textbf{闭合曲线上的线积分}与\textbf{曲面积分}互相转化，是推导分界面衔接条件、建立
“环量 $\Leftrightarrow$ 旋度”场方程的工具。

\corepoint{斯托克斯定理中 $L$ 与 $S$ 的方向必须成右手螺旋；若算出的环量符号与预期相反，八成是绕行方向
或法向取反了。散度定理与斯托克斯定理合起来，构成“微分场方程 $\Leftrightarrow$ 积分场方程”的完整通道。}

\section{无源场和无旋场}

\subsection{无源场}

若在场域中每一点恒有 $\nabla\cdot\vec{A}=0$，称 $\vec{A}$ 为\textbf{无源场}。它有两个重要性质。

\textbf{性质 1（矢量管通量守恒）}：无源场中穿过任一个矢量管的所有截面的通量都相等。
证明要点是取矢量管的一段，其上、下端面 $S_1$、$S_2$ 与侧面 $S_3$ 构成闭合面，由散度定理
$\oiint\vec{A}\cdot\mathrm{d}\vec{S}=0$，而侧面上 $\vec{A}\cdot\vec{e}_n=0$，故穿过 $S_1$ 与
$S_2$ 的通量大小相等、符号相反。在磁场中这一性质称为\textbf{磁通连续性原理}。

\textbf{性质 2（存在矢势）}：任一矢量场旋度再取散度恒为零，即
$\nabla\cdot(\nabla\times\vec{F})\equiv 0$，所以旋度场一定是无源场；反之，无源场 $\vec{A}$ 必存在
另一矢量场 $\vec{F}$ 使
\begin{formulaBox}
\[
\vec{A}=\nabla\times\vec{F}
\]
\end{formulaBox}

$\vec{F}$ 称为 $\vec{A}$ 的\textbf{矢势}（矢量位）。例如恒定磁场中 $\nabla\cdot\vec{B}=0$，
故 $\vec{B}=\nabla\times\vec{A}$，$\vec{A}$ 即磁矢位。

\subsection{无旋场}

若在场域中每一点恒有 $\nabla\times\vec{A}=0$，称 $\vec{A}$ 为\textbf{无旋场}。

\textbf{性质 1}：无旋场沿任意闭合路径的环量等于零，$\oint_L\vec{A}\cdot\mathrm{d}\vec{l}=0$。
因此无旋场中场力做功与路径无关，故又称\textbf{保守场}。静电场 $\oint\vec{E}\cdot\mathrm{d}\vec{l}=0$ 即属此类。

\textbf{性质 2}：无旋场可表示为某一标量函数的梯度场。因为任一标量场的梯度再取旋度恒为零，
即 $\nabla\times(\nabla\varphi)\equiv 0$，所以梯度场一定是无旋场；反之无旋场 $\vec{A}$ 可写成
\begin{formulaBox}
\[
\vec{A}=-\nabla\varphi
\]
\end{formulaBox}

这样的场称\textbf{位势场}，$\varphi$ 称 $\vec{A}$ 的\textbf{位函数}。静电场中 $\varphi$ 就是电位，
满足 $\vec{E}=-\nabla\varphi$。（式中负号来自物理约定：电场指向电位降低的方向。）

\subsection{调和场}

散度与旋度都等于零的矢量场称\textbf{调和场}，它是既无源又无旋的场。设 $\vec{A}=-\nabla\varphi$
且 $\nabla\cdot\vec{A}=0$，则
\begin{formulaBox}
\[
\nabla\cdot(-\nabla\varphi)=-\nabla^2\varphi=0
\quad\Longrightarrow\quad
\nabla^2\varphi=0
\]
\end{formulaBox}

即调和场的位函数必满足\textbf{拉普拉斯方程}；满足拉普拉斯方程且具有二阶连续偏导数的函数称
\textbf{调和函数}。若矢量场仅无旋而不无源，其位函数满足\textbf{泊松方程}
$\nabla^2\varphi=-\rho/\varepsilon_0$。这正是静电场边值问题的方程来源。

\section{亥姆霍兹定理}

\textbf{定理}：在空间有限区域 $V$ 内的任一矢量场 $\vec{A}$，由它的\textbf{散度}、\textbf{旋度}
和\textbf{边界条件}（即包围 $V$ 的闭合面 $S$ 上矢量场的分布）唯一地确定。

等价表述：给定了矢量场的通量源密度 $\nabla\cdot\vec{A}$、涡旋源密度 $\nabla\times\vec{A}$ 以及
场域的边界条件，就可唯一确定该矢量场。

\begin{itemize}
    \item 该定理说明：一个矢量场可以分解为无旋部分与无散部分之和，
    \begin{formulaBox}
    \[
    \vec{A}=-\nabla\varphi+\nabla\times\vec{F}
    \]
    \end{formulaBox}
    前者由散度源（通量源）产生，后者由旋度源（涡旋源）产生。
    \item 它是电磁场理论中“场的唯一性”与“场的分解”的基础：静电场由 $\nabla\cdot\vec{D}=\rho_f$
    与 $\nabla\times\vec{E}=0$ 唯一确定；恒定磁场由 $\nabla\times\vec{H}=\vec{J}$ 与
    $\nabla\cdot\vec{B}=0$ 唯一确定；时变场则由完整的麦克斯韦方程组确定。
    \item 由此也看出为什么静电场可以用标量电位描述（无旋，只剩散度源），而恒定磁场必须用矢量位
    描述（无散，只剩旋度源）。
\end{itemize}

\corepoint{本章五组关系必须烂熟：$\nabla\varphi$（梯度，标量$\to$矢量）、$\nabla\cdot\vec{A}$（散度，矢量$\to$标量）、
$\nabla\times\vec{A}$（旋度，矢量$\to$矢量）、$\nabla^2$（拉普拉斯，标量$\to$标量），以及两条恒等式
$\nabla\times(\nabla\varphi)\equiv 0$、$\nabla\cdot(\nabla\times\vec{A})\equiv 0$。
后面每一章的场方程都由这几个算子写成。}
